\chapter{Expert Validation Survey Instrument}\label{app:survey}

This appendix contains the complete instrument and analysis plan for the
proposed expert validation of the ground truth
(Section~\ref{sec:gtvalidation}). The study was designed but not executed
within the degree project; it is specified here so it can be run without
further design work. A distribution-ready copy is maintained in the project
repository (\texttt{docs/expert\_survey.md}).

\section{Purpose}

Measure (i) how consistently practising cybersecurity professionals assign
the six relationship labels to provision pairs, and (ii) how well their
majority judgement agrees with the project's ground truth. High agreement
validates the ground truth as an evaluation reference; disagreement localises
contestable annotations.

\section{Participants and recruitment}

Target group: professionals with at least two years of experience in
cybersecurity engineering, compliance/GRC, or security assurance, familiar
with at least one of the two standards or an equivalent framework (ISO/IEC
27001, IEC 62443). Target sample: 5--10 respondents recruited through
professional networks and university industry contacts. Participation is
anonymous and uncompensated.

\section{Ethics and data handling}

No personal data is collected beyond role category, years of experience, and
self-rated familiarity. Responses are stored pseudonymously (R1, R2, \ldots).
Participants consent via a checkbox after reading a description of the study
purpose and the intended publication of aggregated results. The design
collects anonymous professional judgements and no sensitive data;
supervisor confirmation is obtained before deployment.

\section{Screening questions}

\begin{enumerate}[label=S\arabic*.]
  \item Current role: security engineer / compliance-GRC / auditor /
  researcher / other (single choice).
  \item Years of professional experience in cybersecurity (number;
  respondents under 2 years excluded from the main analysis).
  \item Familiarity with \standardA{} (Likert 1--5).
  \item Familiarity with \standardB{} or other AI-security guidance
  (Likert 1--5).
  \item Familiarity with compliance or standards mapping as an activity
  (Likert 1--5).
\end{enumerate}

\section{Main task}

Each respondent rates the same 30 provision pairs, drawn from the annotated
sample and stratified over all six labels, in per-respondent random order. For each pair the full
text of both provisions is shown together with the label definitions of
Appendix~\ref{app:labels}, and the respondent answers:

\begin{enumerate}[label=Q\arabic*.]
  \item \textbf{Relationship} --- exactly one of the six labels.
  \item \textbf{Confidence} --- 1 (guessing) to 5 (certain).
  \item \textbf{Justification} --- optional free text.
\end{enumerate}

Estimated completion time is 35--45 minutes. One duplicated pair with
reversed presentation order serves as an attention check; inconsistent
answers flag the respondent for review rather than automatic exclusion.

\section{Post-task questions}

\begin{enumerate}[label=P\arabic*.]
  \item Clarity of the six label definitions (Likert 1--5).
  \item Which pair of labels was hardest to distinguish? (choice among label
  pairs).
  \item Would a tool proposing candidate mappings with these labels help
  your compliance work? (Likert 1--5 plus free text).
\end{enumerate}

\section{Analysis plan}

\begin{enumerate}
  \item \textbf{Inter-expert agreement:} Fleiss' kappa over all respondents,
  on the six-label and the merged five-label scheme, interpreted with the
  Landis--Koch bands~\cite{landis1977kappa}.
  \item \textbf{Expert vs.\ ground truth:} per-item majority vote compared
  with the stored label; percent agreement and Cohen's
  kappa~\cite{cohen1960kappa}. Items without a majority are reported as
  contested.
  \item \textbf{Confidence:} mean confidence per label; point-biserial
  correlation between low confidence and disagreement.
  \item \textbf{Ground-truth revision:} items where the expert majority
  disagrees with the stored label are re-examined; convincing corrections
  enter the ground truth through a documented change log, and the full
  evaluation is re-run to test whether the method ranking of
  Chapter~\ref{ch:results} is stable under the revised reference.
\end{enumerate}

\section{Expected outcomes}

Based on comparable annotation studies, moderate inter-expert agreement
($\kappa \approx 0.4$--$0.6$) is expected on the six-label scheme, higher
after merging, with the overlap/complementarity boundary as the main source
of disagreement. The ground truth is considered validated for its purpose if
majority agreement reaches at least 70\% on the merged scheme.
